\documentclass[Space_Shuttle_Vessel_Manual.tex]{subfiles} 
\begin{document}

\section{SSV and Orbiter Space Flight Simulator}
\begin{multicols*}{2}
\renewcommand{\cfttoctitlefont}{\bf}
\localtableofcontents
\noindent
\\
The section provides information about how to operate Space Shuttle Vessel with the features and limitations of Orbiter Space Flight Simulator.
\newpage


\subsection{SSV User Interface}
This section covers the controls required to interact with the main vessel in SSV, the Orbiter Vehicle. Controls for other SSV vessels are listed in their subsections in Section \ref{sec:ssv-systems};


\subsubsection{Virtual Cockpit controls}
The following paragraphs detail how to operate the virtual cockpit controls with a mouse.

\paragraph{Switches}
Making up the vast majority of the controls in the virtual cockpit, the switches are operated by left-click on the upper-half to move the switch up, or the lower-half to move the switch down. Horizontal switches operate in a similar manner.

\paragraph{Locking Levers}
A locking lever is similar to a switch, except it contains a locking mechanism to avoid accidental operation. When a left-click is initiated, the lever is lifted (unlocked), and then, while maintaining the left-click, a small motion up or down (left or right, for horizontal locking levers) operates the locking lever in the desired direction. Upon release of the left-click, the lever will lock again.

\paragraph{Switch covers}
Left-clicking on a switch cover will open it, allowing operation of the covered control. To close the cover, left-click in it's current location.

\paragraph{Circuit breakers}
Circuit breakers are open or closed by left-clicking on them for 0.2 seconds.

\paragraph{Rotary switches}
Rotary switches come in 2 types, with positions and continuous, but they are operated in the same way. To operate a rotary switch, left-click on the side of the current position in which you want the rotation to happen.

\paragraph{Push buttons}
Push buttons, and push button indicators, are pressed by a left-click.

\paragraph{Selector wheels}
There are 2 types of selector wheels: the original thumb-wheels, and the later push-wheels (which replaced the thumb-wheels, mostly in OV-103 and OV-105). Depending on their use, they can rotate 360º or have a limited range. Thumb-wheels are operated by left-clicking on the upper-half of the wheel, to rotate the wheel up, while left-clicking on the bottom-half of the wheel rotates the wheel down. The push-wheels rotate up or down according to the button that was pressed with a left-click.


\paragraph{DPS keyboards}
The 3 DPS keyboards are operated by left-clicking on the desired key.


\subsubsection{Rotational Hand Controller / Translational Hand Controller}
The regular Orbiter Space Flight Simulator thruster control commands (either keyboard or joystick) are used to simulate the Rotational Hand Controller (RHC) and the Translational Hand Controller (THC). To use the RHC and THC, the associated \textit{FLT CNTLR PWR} switch must be on, as well as the associated (Display Driver Unit) DDU. For the RMS RHC and THC, the RMS needs to be powered.
\\
\\
The RHCs on the real shuttle have a "soft stop" and a "hard stop" (the mechanical limit of movement). Moving the RHC out of detent (up to the soft stop) will command either a constant rotation rate or a pulse of RCS firings to change the rotation rate by a specified amount (depending on whether \textit{DISC RATE} or \textit{PULSE} has been selected). Moving the RHC past the soft stop will result in continuous thruster firings in the appropriate axis. In SSV, a thruster command of <75\% is considered to be within the soft stop; a thruster command of >75\% is treated as RHC deflection beyond the soft stop. When using keyboard controls, the normal keyboard controls are equivalent to full RHC deflection, while holding down the Ctrl key is equivalent to deflection within the soft stop. The THC (and the RMS controls) does not have this idea of a soft stop. When \textit{NORM} is selected in a translation axis, the thrusters will fire continuously if the THC is moved out of detent. When \textit{PULSE} is selected, the thrusters will fire to provide a specified $\Delta$V (the TRAN PLS rate specified on the SPEC 20 DAP CONFIG display). When controlling the RMS, the commanded rotation/translation rates are always directly proportional to the RHC/THC deflection.\\
RHC trim switches are commanded with the Shift+Arrow Key.

\begin{figure}[H]
  \includegraphics[width=0.99\hsize]{RHC_THC_keys.png}
  \caption{RHC/THC key mapping}
  \label{fig:RHC_THC_keys}
\end{figure}

\subsubsection{Speedbrake/Thrust Controller}
The Speedbrake/Thrust Controller (SBTC) controls both SSME throttling during ascent and the speedbrake setting during entry. In Orbiter, the SBTC is controlled in 5\% intervals, in the forward direction by the numpad \textit{Subtract} key, and in the aft direction via the numpad \textit{Add} key. The takeover switch, used to initiate manual control of the SSME throttle or speedbrake settings, is simulated by the \textit{Minus} key. The appropriate \textit{FLT CNTLR PWR} switch must be on, as well as the appropriate DDU, for the SBTC to be active. Each SBTC is animated so the user can tell what setting is being commanded. During ascent, a full aft SBTC position corresponds to MPL (Minimum Power Level), usually 65 or 67\% SSME throttle; full forward SBTC position corresponds to the current maximum SSME throttle, a value ranging from 100 to 109\%, depending on mission or abort case. During landing, full aft SBTC position corresponds to the speedbrake being \textbf{fully open}; full forward SBTC position corresponds to the speedbrake being commanded \textbf{fully closed}.

\paragraph{Ascent}
During ascent, SSME throttling is usually controlled by autopilot; in this case, the \textit{AUTO} portion of the \textit{SPD BK/THROT} PBIs on Panel F2 \& Panel F4 is lit, and \textbf{THRTL: Auto} is displayed in the A/E PFD display. To takeover manual control of the SSME throttle command, press the SBTC takeover switch and move the SBTC to match the current SSME auto command (displayed in the Ascent Traj displays). When the SBTC takeover switch is pressed both \textit{AUTO} PBIs will go out, indicating a manual takeover is in progress, but not completed. When the SBTC-commanded SSME throttle setting matches the auto command within 4\%, the PLT \textit{SPD BK/THROT MAN} PBI will be lit and \textbf{THRTL: Man} appears in the A/E PFD display, indicating that the SSME throttle is now under manual control. To return to auto SSME throttle control, press either \textit{SPD BK/THROT MAN} PBI. A manual MECO can be commanded by pressing the NUMPAD * key (in real life, this is done by simultaneously pressing all 3 \textit{MAIN ENGINE SHUT DOWN} push buttons on Panel C3; this is not possible in Orbiter).

\paragraph{Entry}
\label{sec:entry}
The speedbrake is usually controlled automatically throughout entry, with the \textit{AUTO} portion of the \textit{SPD BK/THROT} PBIs on Panel F2 \& Panel F4 is lit, and \textbf{SB: Auto} is displayed in the A/E PFD display. To take over manual control of the speedbrake press the SBTC takeover switch; the speedbrake will immediately move to the position commanded by the SBTC, the \textit{AUTO} portion of the \textit{SPD BK/THROT} PBIs will go out and the \textit{MAN} PBI will be lit on either the CDR or PLT position (depending on what SBTC takeover switch was last pressed). Pressing the either \textit{SPD BK/THROT} PBI, will put the speedbrake into \textit{AUTO} mode again.

\subsubsection{Rudder Pedal Transducer Assembly}
The Rudder Pedal Transducer Assembly (RPTA) allows manual control of the rudder during the later part of reentry, as well as the Nose Wheel Steering during rollout. The RPTA also contains the brake pedals, which in addition to braking provide another means of lateral control during rollout. Although the rudder is automatically controlled, manual control is available when the R/Y channel is in CSS. The appropriate \textit{FLT CNTLR PWR} switch must be on, as well as the appropriate DDU, for the RPTA to be active.


\begin{figure}[H]
  \includegraphics[width=0.99\hsize]{SBTC_RPTA_keys.png}
  \caption{SBTC/RPTA key mapping}
  \label{fig:SBTC_RPTA_keys}
\end{figure}


\subsubsection{Additional keyboard controls}
In addition to the controls described above, there are a few more keys needed to operate the vehicle.\\
\\
\textbf{General}\\
Ctrl+A - toggle between CDR/PLT/AFD RHC and THC, and RMS RHC and THC\\
Ctrl+G - arm landing gear\\
G - deploy landing gear\\
\\
\\
\textbf{RMS}\\
Ctrl+Enter - grapple\\
Ctrl+Backspace - release\\
Ctrl+O - toggle between Coarse and Vernier rates
\newpage


\subsection{Navigating the Virtual Cockpit}
\label{sec:nav-vc}
SSV has several Virtual Cockpit (VC) positions, allowing the user to navigate between the main locations and station of the crew module. Table \ref{tab:VC_navigation} lists all the available VC positions, as well as neighbour positions accessed with the standard Ctrl+Arrow key combination. Each position also has lean positions (Alt Gr+Arrow key) which allow better positioning to reach a specific panel or window.\\
For assistance in navigating the VC, the name of the position is shown for a few seconds at the top of the screen during the simulation.
\\

\end{multicols*}
\begin{table}[H]
\begin{threeparttable}
  \centering
  \begin{tabular}{l|p{2.88cm} p{2.88cm} p{2.88cm} p{2.88cm} }
	\textbf{Cockpit View} & \textbf{Left} & \textbf{Right} & \textbf{Up} & \textbf{Down} \\
	\hline\rule{0pt}{2ex}
	Commander Seat & CDR - L4 & Pilot Seat & PLB Camera A \textit{or} ODS Camera\tnote{c} & MS Seat \\
	\hline\rule{0pt}{2ex}
	Pilot Seat & Commander Seat & Pilot - R4 & PLB Camera D \textit{or} ODS Camera\tnote{c} & MS2/FE Seat \\
	\hline\rule{0pt}{2ex}
	CDR - L4 & Port Workstation & Commander Seat & PLB Camera D \textit{or} ODS Camera\tnote{c} & MS Seat \\
	\hline\rule{0pt}{2ex}
	Pilot - R4 & Pilot Seat & Stbd Work Station & PLB Camera D \textit{or} ODS Camera\tnote{c} & MS2/FE Seat \\
	\hline\rule{0pt}{2ex}
	MS Seat & Port Workstation & MS2/FE Seat & Commander Seat & PLB Camera A \textit{or} ODS Camera\tnote{c}\\
	\hline\rule{0pt}{2ex}
	MS2/FE Seat & MS Seat & Stbd Work Station & Pilot Seat & PLB Camera A \textit{or} ODS Camera\tnote{c}\\
	\hline\rule{0pt}{2ex}
	Port Work Station & RMS Work Station & Commander Seat & PLB Camera A \textit{or} ODS Camera\tnote{c} & Mid Deck\\
	\hline\rule{0pt}{2ex}
	Stbd Work Station & Pilot Seat & Aft Pilot Station & PLB Camera D \textit{or} ODS Camera\tnote{c} & Aft Work Station\\
	\hline\rule{0pt}{2ex}
	Aft Work Station & Stbd Work Station & Port Workstation & RMS Work Station & MS Seat\\
	\hline\rule{0pt}{2ex}
	Aft Pilot Station & Stbd Work Station & RMS Work Station & PLB Camera D \textit{or} ODS Camera\tnote{c} & Aft Work Station\\
	\hline\rule{0pt}{2ex}
	RMS Work Station & Aft Pilot Station & Port Work Station & PLB Camera A \textit{or} ODS Camera\tnote{c} & Aft Work Station\\
	\hline\rule{0pt}{2ex}
	RMS EE\tnote{a} & RMS Elbow & - & - & RMS Work Station\\
	\hline\rule{0pt}{2ex}
	RMS Elbow\tnote{a} & - & RMS EE & - & RMS Work Station\\
	\hline\rule{0pt}{2ex}
	PLB Camera A & PLB Camera D & PLB Camera B & RMS EE\tnote{a,c} & RMS Work Station \textit{or} ODS Camera\tnote{c}\\
	\hline\rule{0pt}{2ex}
	PLB Camera B & PLB Camera A & PLB Camera C & RMS EE\tnote{a,c} & RMS Work Station \textit{or} ODS Camera\tnote{c}\\
	\hline\rule{0pt}{2ex}
	PLB Camera C & PLB Camera B & PLB Camera D & RMS EE\tnote{a,c} & Aft Pilot Station \textit{or} ODS Camera\tnote{c}\\
	\hline\rule{0pt}{2ex}
	PLB Camera D & PLB Camera C & PLB Camera A & RMS EE\tnote{a,c} & Aft Pilot Station \textit{or} ODS Camera\tnote{c}\\
	\hline\rule{0pt}{2ex}
	ODS Camera\tnote{c} & - & - & PLB Camera A & Aft Pilot Station\\
	\hline\rule{0pt}{2ex}
	Mid Deck & - & - & Port Work Station & External Airlock\tnote{b}\\
	\hline\rule{0pt}{2ex}
	External Airlock\tnote{b} & - & - & Mid Deck & ODS Camera\tnote{c}\\
  \end{tabular}
	\begin{tablenotes}
		\item[a] only when RMS is installed
		\item[b] only when External Airlock is installed
		\item[c] only when ODS is installed
	\end{tablenotes}
	\end{threeparttable}
  \caption{VC navigation}
  \label{tab:VC_navigation}
\end{table}
\newpage


\subsection{Atmospheric Wind}
\begin{multicols*}{2}
\noindent
Since the 2016 version, Orbiter Space Flight Simulator implements atmospheric winds.\\
Unlike the real world, there is no weather forecast in Orbiter, so during landing it is easy to fly the vehicle into a wind situation that exceeds its (narrow) operating capabilities. For this reason, if atmospheric wind effects are enabled, intense monitoring of the vehicle performance during landing is required to ensure the vehicle reaches the runway.\\
\end{multicols*}
\newpage


% not upstream: Flight recorder section (SSV Linux port, flight recorder support)
\subsection{Flight recorder}
\label{sec:flight-recorder}
\begin{table}[b!]
  \centering
  \begin{tabularx}{470pt}{>{\centering\arraybackslash}p{75pt} | >{\centering\arraybackslash}p{100pt} | >{\centering\arraybackslash}p{150pt} | >{\centering\arraybackslash}p{95pt}}
    \textbf{Event} & \textbf{Text} & \textbf{Written by} & \textbf{On playback} \\
    \hline
	SSV\_NOMINAL & format (1) & every SSV vessel, at the start of a recording & nominal values, user input locked\\
	\hline\rule{0pt}{2ex}
	SSV\_VC & "panel" "component" state & VC switches, breakers, rotaries, wheels, pushbuttons, keyboards, MDUs, CTVM & the component is set\\
	\hline\rule{0pt}{2ex}
	SSV\_CTRL & numbers & OV hand controllers, RPTA, SBTC, brakes and RMS controllers; crawler and OTS drive; Centaur and IUS manual gimbal & the controls are set\\
	\hline\rule{0pt}{2ex}
	SSV\_CMD & command and arguments & keyboard commands, pad dialogs, LCC MFD & the vessel runs the command\\
	\hline\rule{0pt}{2ex}
	SSV\_EVENT, SSV\_CW, SSV\_MA & text & mission events, C\&W messages, master alarm & nothing (notes)
	\end{tabularx}
  \caption{Flight recorder events}
  \label{tab:FlightRecorderEvents}
\end{table}

\begin{multicols*}{2}
\noindent
SSV uses the Orbiter flight recorder (Ctrl+C) to record a mission and play it back.\\
\\
\textbf{What is recorded}\\
Orbiter records the trajectory, attitude, thrust levels, attachments and its own events. SSV adds the crew's actions: the VC controls, the hand controllers, pedals and brakes, the keyboard commands, the pad dialogs, the crawler and OTS controls and the LCC MFD. Mission events, C\&W messages and the master alarm are written as notes, which can be read in the vessel's .atc file in "<Orbiter installation>\textbackslash Flights".\\
\\
\textbf{Playback}\\
During a playback the SSV systems run live on the recorded trajectory and inputs. Random variation (sensor errors, APU and SRB noise, ET boil-off, MPS helium supply pressure, SSME helium usage, MLP igniter time, Xenon light aim) is set to its nominal value, so nothing random can cause a hold or a C\&W that the recording did not have. A playback repeats the inputs and states of the recording, not every frame of it.\\
\\
\textbf{Input lock}\\
While a vessel plays back, its VC, keyboard commands, dialogs and MFD commands ignore the user, as the recorded crew is in control. Camera and view keys still work. The lock ends with the playback.\\
\\
\textbf{Limits}\\
Recorded events play back one or two frames after their recorded time. Vessels created during a recording are not recorded (Orbiter's limit). A recording made with an SSV version without this support plays back without the nominal values and the lock.\\
\\
\textbf{Events}\\
Table \ref{tab:FlightRecorderEvents} lists the events SSV writes, for add-on authors who read or write .atc files. Each line is "<time> <event> <text>". Names in the text are in double quotes, as in the scenario file.\\
\\
\textbf{SSV\_CTRL values}\\
OV: 35 values, RHC (left, right, aft: pitch, roll, yaw), THC (left, aft: X, Y, Z), trims (left, right: pitch, roll), RPTA, SBTC position and takeover (left, right), brakes (CDR, PLT: left, right), RMS RHC, RMS THC, RMS single joint drive.\\
Crawler: target speed (front, rear cab), steering (front, rear cab), brake, jack height index.\\
OTS: front steering, rear steering, drive.\\
Centaur: pitch, yaw and roll rate commands. IUS: pitch and yaw rate commands.\\
\\
\textbf{SSV\_CMD commands}\\
OV: GEAR ARM/DOWN, SSMEPB, RMS 0/1, RMSEE GRAPPLE/RELEASE/OFF, RMSRATE 0/1, PLBCAM <left> <right> <up> <down>.\\
LCC: RESUME, CUTOFF, HOLD <0730/0500/0400/0255/0157/0031> 0/1, FIREX <1-3> 0/1.\\
LC-39 and SLC-6: OAA EXTEND/RETRACT/HALT, GVA EXTEND/RETRACT/HALT, HOOD LOWER/RAISE/HALT, GVAHOOD EXTEND/RETRACT, ETVAS ATTACH/DETACH, IAA DEPLOY/HALT/RETRACT.\\
LC-39: RSS OPEN/STOP/CLOSE, RSS SET OPEN/CLOSE (Ctrl+K), RSSOWP OPEN/CLOSE, FSSOWP OPEN/CLOSE, RBUS EXTEND/RETRACT, SRBSFD EAST/WEST EXTEND/RETRACT.\\
SLC-6: PCR, SAB and MST EXTEND/HALT/RETRACT, SABDOOR OPEN/HALT/CLOSE, EAA DEPLOY/HALT/RETRACT, SDWW LOWER/RAISE.\\
Xenon Lights: POWER 0/1.\\
Centaur: PLSEP, PL <1-8>, RL10 START/STOP.\\
IUS: PLSEP, STAGE, PL <1-8>, SRM START 1/2.
\end{multicols*}
\newpage


\subsection{Options}
\begin{table}[b!]
  \centering
  \begin{tabularx}{470pt}{c | >{\centering\arraybackslash}p{40pt} | >{\centering\arraybackslash}p{40pt} | >{\centering\arraybackslash}p{240pt}}
    \textbf{Parameter} & \textbf{Type} & \textbf{Default} & \textbf{Description} \\
    \hline
	EIUDataRecorder & bool & FALSE & Enables output of SSME controller data to file "<Orbiter installation>\textbackslash EIU\_data\_ch<id>.txt"\\
	\hline\rule{0pt}{2ex}
	PositionLabelTime & integer & 3 & Controls duration of vc position label output, [0;60] seconds\\
	\hline\rule{0pt}{2ex}
	AutoActionLandingGear & bool & TRUE &Enables automatic Landing Gear arm and deploy\\
	\hline\rule{0pt}{2ex}
	AutoActionDragChute & bool & TRUE &Enables automatic Drag Chute deploy and jettison\\
	\hline\rule{0pt}{2ex}
	RHCRate & double & 0.75 & Controls keyboard input rate to RHC (1/x seconds from detent to full deflection)\\
	\hline\rule{0pt}{2ex}
	RPTARate & double & 0.75 & Controls keyboard input rate to RPTA (1/x seconds from detent to full deflection)\\
	\hline\rule{0pt}{2ex}
	BrakeRate & double & 0.75 & Controls keyboard input rate to Brake (1/x seconds from detent to full deflection)\\
	\hline\rule{0pt}{2ex}
	Alarm1x & bool & TRUE & Enables time acceleration being set to 1x when a C\&W alarm is activated
	\end{tabularx}
  \caption{Option Parameter List}
  \label{tab:OptionParameterList}
\end{table}

\begin{multicols*}{2}
\noindent
The file "<Orbiter installation>\textbackslash Config\textbackslash Vessels\textbackslash SSV\_OV.cfg" contains parameters to control some features of the Orbiter Vehicle vessel.\\

\WARNING{The "OV.cfg" file will be overwritten during the installation of a new SSV version.}
\\
\\
\noindent
Table \ref{tab:OptionParameterList} lists all the available option parameters.

\end{multicols*}

\end{document}